\section{What Makes a Visualization Memorable?}

Der Artikel welcher im Journal IEEE Transactions on Visualization and Computer Graphics 2013 veröffentlicht wurde, wurde ausgewählt, weil Daten ihren Nutzen erst dann voll entfalten, wenn auch ihr Kontext nachvollziehbar vermittelt wird. Da Menschen Informationen in hohem Maß visuell aufnehmen und verarbeiten, stellt sich für unser Projekt unmittelbar die Frage, wie Daten so visualisiert werden können, dass sie für unterschiedliche Nutzer:innen verständlich und zugänglich sind. Dies ist insbesondere deshalb relevant, weil unser Vorhaben darauf abzielt, Open Data auch für Laien nutzbar zu machen. Verständliche und zugleich einprägsame Visualisierungen sind dafür eine zentrale Voraussetzung.

Der Artikel \textit{What Makes a Visualization Memorable?} untersucht in diesem Zusammenhang, welche Eigenschaften dazu beitragen, dass Visualisierungen im Gedächtnis bleiben und welche Rolle dies für die Wahrnehmung und Rezeption von Daten spielt. Untersucht wurde diese Fragestellung in einer Studie mit $n = 261$ Teilnehmenden. \cite{Visualization}

\smallskip
\noindent
\paragraph{Take-Aways}
\begin{enumerate}
    \item \textbf{Die Eigenschaft der Memorierbarkeit}

    Der Artikel zeigt, dass die Merkf\"ahigkeit von Visualisierungen als eigene Eigenschaft verstanden werden kann. Die Studie zeigt, dass Personen relativ konsistent darin sind, welche Visualisierungen sie als merkbar oder weniger merkbar wahrnehmen (vgl.\ Abstract sowie Section~1: Introduction und Section~8: Conclusions and Future Work in \cite{Visualization}).

    $\rightarrow$ F\"ur unser Projekt bedeutet das, dass wir nicht nur auf Korrektheit und Usability achten sollten, sondern auch darauf, ob Inhalte \"uberhaupt im Ged\"achtnis bleiben.

    \item \textbf{Farbe und erkennbare Objekte erh\"ohen die Merkf\"ahigkeit}

    Ein zentrales Ergebnis der Studie ist, dass farbige Visualisierungen sowie Darstellungen mit menschlich erkennbaren Objekten oder Piktogrammen signifikant besser erinnert werden. Solche visuellen Anker machen Visualisierungen auff\"alliger und erleichtern die Wiedererkennung (vgl.\ Abstract sowie Section~8: Conclusions and Future Work in \cite{Visualization}).

    $\rightarrow$ F\"ur unsere App hei\ss t das, dass Farben und bildhafte Elemente bewusst eingesetzt werden sollten, um Daten zug\"anglicher und einpr\"agsamer zu machen.

    \item \textbf{Ungew\"ohnliche Visualisierungstypen bleiben eher im Ged\"achtnis}

    Die Studie zeigt au\ss erdem, dass ungew\"ohnlichere Darstellungsformen wie Diagramme, Grid-/Matrix-Visualisierungen oder Trees/Networks besser erinnert werden als klassische Standarddiagramme wie Balken-, Linien- oder Punktdiagramme. Neuartige und unerwartete Formen scheinen damit st\"arker im Ged\"achtnis zu bleiben (vgl.\ Abstract sowie Section~8: Conclusions and Future Work in \cite{Visualization}).

    $\rightarrow$ F\"ur unser Projekt ist das ein Hinweis, dass wir nicht nur auf Standarddiagramme setzen sollten. Eine innovative App kann sich gerade durch passend gew\"ahlte, weniger konventionelle Darstellungsformen von klassischen Open-Data-Portalen abheben.
\end{enumerate}
 
\smallskip
\noindent
Zusammenfassend zeigt der Artikel, dass einprägsame Visualisierungen häufig farbig, bildhaft, visuell markant und weniger standardisiert sind. Gleichzeitig macht die Studie deutlich, dass Merkfähigkeit nicht mit analytischer Qualität gleichgesetzt werden darf. Für unser Projekt ist das besonders hilfreich, weil wir keine reine Analyseplattform bauen, sondern eine mobile App, die Open Data für eine breite Öffentlichkeit zugänglicher und interessanter machen soll. Daraus lässt sich ableiten, dass unsere Visualisierungen nicht nur korrekt und verständlich, sondern auch visuell erinnerbar gestaltet sein sollten.


